% Editable chapter, migrated from the canonical Phase 9B report.
\section{Analyse exploratoire et stationnarité}
\label{sec:eda}

\subsection{Dispersion et dépendance des rendements}



\historicalfigure{log_returns.png}{Log-rendements des trois actifs}{Variations quotidiennes des prix ajustés. Source : \sourcepath{outputs/figures/phase2/log_returns.png}. Périmètre : descriptif complet, 2018-01-02--2026-10-02, TEST marqué. Unité : pourcentage logarithmique quotidien. Les amplitudes regroupées suggèrent une dynamique de variance ; elles n’identifient pas une cause économique.}{0.32}



Les rendements de la figure~\ref{fig:log_returns.png} oscillent autour de niveaux bien plus faibles que leur dispersion. Certaines séquences concentrent de grands mouvements, puis laissent place à des périodes plus calmes. Ce volatility clustering motive des diagnostics de variance ; le graphique seul ne constitue pas un test ARCH. Les statistiques utilisées pour les choix sont celles de TRAIN, même lorsque le contexte visuel couvre l’ensemble de la période.



% Documentary display only; edit preparation definitions for a refresh.
% Source CSVs: outputs/tables/phase2_descriptive_stats.csv
\begin{table}[!htbp]
\centering
\caption{Distribution des log-rendements sur TRAIN.}
\label{tab:descriptive}
\small
\setlength{\tabcolsep}{4pt}
\renewcommand{\arraystretch}{1.16}
\begin{tabularx}{\linewidth}{Xrrrr}
\toprule
\textbf{Actif} & \textbf{n} & \textbf{Moyenne \%} & \textbf{Écart-type \%} & \textbf{Kurtosis exc.} \\
\midrule
NVDA & 1979 & 0,185 & 3,244 & 4,698 \\
AMD & 1979 & 0,157 & 3,451 & 3,571 \\
SOXX & 1979 & 0,085 & 2,219 & 5,126 \\
\bottomrule
\end{tabularx}
\par\smallskip
\begin{minipage}{\linewidth}\footnotesize Source : \sourcepath{outputs/tables/phase2_descriptive_stats.csv}, filtre \sourcepath{scope=train}, \sourcepath{series=Log_Return}. Moyenne et écart-type en pourcentage quotidien ; kurtosis excédentaire sans unité ; arrondi à trois décimales.
\end{minipage}
\end{table}




Le tableau~\ref{tab:descriptive} sépare deux dimensions du risque. AMD a la dispersion la plus élevée, tandis que SOXX a la kurtosis excédentaire la plus élevée. Une kurtosis élevée décrit des queues épaisses relatives à la forme normale et ne signifie pas nécessairement une volatilité globale supérieure. Les corrélations de Pearson contemporaines sur TRAIN sont positives : NVDA/AMD 0,685, NVDA/SOXX 0,819 et AMD/SOXX 0,745. Ces associations motivent l’étude multivariée, sans démontrer de causalité.



Les seuils IQR calculés sur TRAIN repèrent 74 rendements atypiques pour NVDA, 89 pour AMD et 74 pour SOXX. Ils sont tous conservés : une observation atypique peut être une variation de marché réelle. La supprimer automatiquement modifierait précisément le risque que l’étude cherche à décrire.

La figure~\ref{fig:rolling_correlations_60d.png} en annexe B estime les corrélations contemporaines sur des fenêtres rétrospectives de 60 séances TRAIN. Leurs moyennes valent 0,669 (NVDA/AMD), 0,796 (NVDA/SOXX) et 0,738 (AMD/SOXX) ; les plages respectives sont [0,154 ; 0,936], [0,506 ; 0,956] et [0,348 ; 0,943]. Une corrélation globale masque donc des variations temporelles. Les fenêtres se chevauchent et leurs estimations ne sont pas indépendantes ; elles ne démontrent ni causalité économique ni gain prédictif du VAR. Source : \sourcepath{outputs/tables/phase2_rolling_correlation_summary.csv}.



Sources : \sourcepath{outputs/tables/phase2_return_correlation_train.csv}, \sourcepath{outputs/tables/phase2_outlier_summary.csv}.





\subsection{ACF PACF et tests complémentaires}



La forte persistance des niveaux apparaît dans l’ACF des prix ajustés, reproduite en annexe A. Après différenciation logarithmique, la dépendance linéaire diminue fortement. L’ACF mesure l’association avec chaque retard ; la PACF isole la contribution d’un retard après les retards intermédiaires. Des pics isolés au-delà de bandes pointwise ne suffisent pas à fixer un ordre : ces bandes ne contrôlent pas simultanément tous les retards inspectés.



\historicalfigure{nvda_log_return_acf_pacf_train.png}{ACF et PACF des rendements NVDA}{Dépendance après différence logarithmique. Source : \sourcepath{outputs/figures/phase2/nvda_log_return_acf_pacf_train.png}. Périmètre : TRAIN, 2018-01-02--2025-11-14, 1979 rendements valides. Unité : corrélation sans dimension ; retards en séances, bandes ponctuelles à 95 \%.}{0.32}



ADF teste une racine unitaire ; KPSS teste la stationnarité autour des termes déterministes spécifiés~\cite{R4,R5}. Les niveaux utilisent constante et tendance ; les rendements une constante. L’ADF choisit ses retards par AIC et KPSS par sa règle automatique. Leur lecture conjointe est plus informative qu’un test isolé.



% Documentary display only; edit preparation definitions for a refresh.
% Source CSVs: outputs/tables/phase2_stationarity_summary.csv
\begin{table}[!htbp]
\centering
\caption{Synthèse ADF/KPSS.}
\label{tab:stationarity}
\small
\setlength{\tabcolsep}{4pt}
\renewcommand{\arraystretch}{1.16}
\begin{tabularx}{\linewidth}{XXXXX}
\toprule
\textbf{Série} & \textbf{Actifs} & \textbf{ADF} & \textbf{KPSS} & \textbf{Lecture conjointe} \\
\midrule
Adj\_Close & NVDA, AMD, SOXX & Non rejet & Rejet & Non-stationnarité soutenue \\
Log\_Price & NVDA, AMD, SOXX & Non rejet & Rejet & Non-stationnarité soutenue \\
Log\_Return & NVDA, AMD, SOXX & Rejet & Non rejet & Stationnarité soutenue \\
\bottomrule
\end{tabularx}
\par\smallskip
\begin{minipage}{\linewidth}\footnotesize Source : \sourcepath{outputs/tables/phase2_stationarity_summary.csv}. Périmètre : TRAIN ; trois actifs. Résultats catégoriels sans unité ; aucune nouvelle p-value calculée.
\end{minipage}
\end{table}




Pour NVDA log-prix, ADF donne p=0,461103 et KPSS une borne publiée de 0,01, avec p réelle inférieure. Pour son log-rendement, ADF donne environ 7,91 \(\times\) 10\textsuperscript{-27} ; KPSS publie 0,10, avec p réelle supérieure. Les tests soutiennent d=1 sans prouver un processus exactement I(1). Les bornes KPSS ne sont pas des probabilités exactes ; les zéros numériques ne désignent pas une probabilité mathématiquement nulle.



\subsection{Lecture et limites de la décomposition STL}

La figure~\ref{fig:nvda_log_price_stl_train.png} en annexe A décompose le log-prix NVDA TRAIN selon \(y_t=T_t+S_t+R_t\), avec une période imposée de cinq séances et une estimation robuste. La série observée combine une hausse de long terme et des replis, notamment fin 2018 et en 2022. La tendance lissée restitue ces mouvements dominants en atténuant les fluctuations rapides. La composante saisonnière oscille autour de zéro avec une amplitude faible relativement aux niveaux et variable dans le temps. Les résidus, différence entre observation, tendance et saisonnalité, présentent des pics et une dispersion irrégulière ; ils ne sont pas démontrés comme bruit blanc par cette figure.

Cette lecture concorde avec les diagnostics Phase 2 : ADF/KPSS soutiennent la différenciation des log-prix, tandis que l’ACF des rendements au carré et les extrêmes conservés motivent l’étude de la variance. La robustesse du lissage ne supprime ni les ruptures ni les chocs. Ses filtres bilatéraux utilisent des observations ultérieures à chaque date à l’intérieur de TRAIN : les composantes restent descriptives et ne servent pas de features causales. Cinq séances ne coïncident pas toujours avec une semaine calendaire en présence de jours fériés. La période est imposée, les différences par jour de semaine sont descriptives et aucun test Phase 2 n’établit un cycle stable exploitable. STL ne suffit donc pas à justifier SARIMA ; la décision gelée est reprise au chapitre~\ref{sec:arima}.



Sources : \sourcepath{outputs/tables/phase2_stationarity_tests.csv}, \sourcepath{outputs/tables/phase2_weekday_returns.csv}, \sourcepath{docs/phase2_eda_stationarity.md}.
